\documentclass[10pt]{article}

\usepackage[utf8]{inputenc}

\usepackage[T1]{fontenc}

\usepackage{amsmath}

\usepackage{amsfonts}

\usepackage{amssymb}

\usepackage[version=4]{mhchem}

\usepackage{stmaryrd}

\usepackage{graphicx}

\usepackage[export]{adjustbox}

\graphicspath{ {./images/} }



\begin{document}

нуля только коэффициенты с номерами $n_{m}$, образующими лакунарную последовательность, т. е. такими, что $n_{m+1} / n_{m} \geqslant \lambda>1$. Другой пример специальных рядов ряды с монотонными коэффициентами.



Все сказанное выше относилось к Ф. р. вида (2). Для Ф. р. по переставленной тригонометрич. системе нек-рые свойства Ф. р. по тригонометрич. системе, взятой в обычном порядке, не имеют места. Напр., существует такая непрерывная функция, тто ее Ф. р. после нек-рой перестановки расходится почти всюду (cM. [12] - [15]).



Теория Ф. р. для функций многих переменных разработана в меньшей степени. Часть многомерных результатов аналогична одномерным. Но имеются и существенные отличия.



Пусть $x=\left(x_{1}, \ldots, x_{N}\right)$-точка $N$-мерного пространства, $k=\left(k_{1}, \ldots, k_{N}\right)$ есть $N$-мерный вектор с целочисленпыми координатами и $(k, x)=k_{1} x_{1}+\ldots+k_{N} x_{N}$. Для функции $f(x)$, имеющей период $2 \pi$ по каждой переменной и интегрируемой по Лебегу на $N$-мерном губе $[0,2 \pi]^{N}$, Ф. р. по тригонометрич. системе наз. рял



\[

\sum_{k} c_{k} e^{i(k, x)}

\]



где суммирование ведется по всем $k$ и



\[

c_{k}=\frac{1}{(2 \pi)^{N}} \int_{0}^{2 \pi} \cdots \int_{0}^{2 \pi} f(x) e^{-i(k, x)} d x

\]



\begin{itemize}

  \item к о ә фф и ц и е т ты Ф у р в е функции $f . \Phi . \mathrm{p}$. (8) записан в комплексной форме. Запись его в тригонометрич. форме как ряда по произведениям косинусов и синусов кратных дуг более громоздка.

\end{itemize}



Возможны различные определения частных сумм ряда (8), напр. частные суммы по прямоугольникам



по кругам



\begin{center}

\includegraphics[max width=\textwidth]{2023_07_09_991621900b1f6760476ag-1(1)}

\end{center}



\[

\sum_{|k| \leqslant n} c_{k} e^{i(k, x)}

\]



где $n$-радиус и $|k|=\sqrt{k_{1}^{2}+\ldots+k_{N}^{2}}$.



Для иредставления функциї круговые частные суммы (9) менее пригодны, тем их средние Рисса



\[

\sum_{|k| \leqslant n}\left(1-\frac{|k|}{n}\right)^{\alpha} c_{k} e^{i(k, x)} .

\]



Для средних Рисса порядка $\alpha \geqslant \frac{N-1}{2}$ Ф. р. Функций из $L^{2}$ справедлив принцип локализации, а для меньших $\alpha$ это не так (С. Бохнер, S. Bochner, 1936). Среднис Рисса круговых частных сумм критич. порядка $\alpha=\frac{N-1}{2}$ играют существенную роль и в других вопросах Ф. p. функций многих переменных.



Существует непрерывная функция двух переменных, Ф. р. к-рой не сходится по прямоугольникам ни в одной внутренней точке квадрата $[0,2 \pi]^{2}$ (см. [16]).



Нек-рые результаты, относящиеся к Ф. р. по тригонометрич. системе, допускают значительные обобщения, напр. могут быть соответствующим образом перенесены на спектральные разложения, отвечаюощие самосопряженным эллиптическим дифференциальным операторам.



Лuт.. [1] Б а р и Н. К., Тригонометрические ряды, М., англ., 2 изд., т. 1-2, М., 1965; [3] X а p д и Г. Х., Р ог оз и н с к и й В. В., Ряды Фурье, пер. с англ., М., 1959: [4] ЈI узи и Н. Н., Интеграл и тригонометрический ряд, М.- Јі., $1951,[5] \mathrm{L}$ e b e s g u e H., Leçons sur les séries trigonométriques, P., 1906; [6] П а п л а у с к а с А. Б., Тригонометрические ряды от Әйлера до Јебега, М., 1966 ; [7] у л ь я н о в П. Ј., «Успехи матем. наук», 1964, т. 19 , в. 1 , с. $3-69 ;$ [8] А ли и1976, T. 31 , в. 6, c. $28-83$, [9] S a l e m R., "Duke Math. J.",



\begin{center}

\includegraphics[max width=\textwidth]{2023_07_09_991621900b1f6760476ag-1}

\end{center}



C. Б., "Докл. AH CCCP», 1955, T. 102, c.. 37-40; [12] K o 1-

m o g o r of f A., M e n s c h of f D., "Math. Z.", $1927, \mathrm{Bd} 26$, S. $432-41$; [13] i a h or s k i 7 , "C. r. Acad. sci.", 1960, t 251 , p. 501-03; [14] У л ь я н о в П Л., "Успехи матем. паук", "Докл. АН СССР», 1961, т. 141, c. 28-31, [16] F e f f e rm a n C., «Bull. Amer. Math. Soc.», 1971, v. 77, p. 191-95.



ФУРБЕ ЧиСлО - один из критериев подобия нестационарных тепловых процессов. Характеризует соотношение между скоростыю изменения тепловых условий в окружающей среде и скоростью перестройки поля температуры внутри рассматриваемой системы (тела), к-рый зависит от размеров тела и коәффициента его теплопроводности. Ф. ч. $\boldsymbol{F}_{0}=a t_{0^{\prime}} l^{2}$, где $a=\lambda / \rho c-$ коэффициент температуропроводности, $\lambda-$ коәффициент теплопроводности, $\rho$ - плотность, $c$ - удельная теплоемкость, $l$ - характерный линейный размер тела, $t_{0}$ - характерное время шзменения внешних условий.



Название по имени Ж. Фурье (Ј. Fourier). По материалам одноименной статъи БCЭ-з.



ФУРБЕ - БЕССЕЛЯ ИНТЕГРАЛ, И н т е г р а л Г а н к е л я,- аналог Фурье интеграла для Бесселя буғкций, имегщий вид



\[

f(x)=\int_{0}^{\infty} \lambda J_{v}(\lambda x) d \lambda \int_{0}^{\infty} y J_{v}(\lambda y) f(y) d y .

\]



Формула (*) может быть получена из Фурье-Бесселя ряда для интервала $(0, l)$ переходом к пределу при $l \rightarrow$ $+\infty$. Г. Ганкель (H. Hankel, 1875) установил теорему: если функция $f(x)$ кусочно непрерывная и имеет ограниченную вариацию на любом интервале $0<x<l$, интеграл



\[

\int_{0}^{\infty} \sqrt{x}|f(x)| d x

\]



сходится, то формула (\textit{) справедлива при $v>-1 / 2$ во всех точках непрерывности $f(x), 0<x<+\infty$. В точках разрыва $x_{0}, 0<x_{0}<+\infty$, правая часть формулы (}) равна $\left[f\left(x_{0}-0\right)+f\left(x_{0}+0\right)\right] / 2$, при $x_{0}=0$ она дает $f(0+) / 2$.



Аналоги Ф. - Б. и. (*) для цилиндрит. функций $Z_{v}(x)$ пругих типов также справедливы, но пределы интегралов должны быть соответственно изменены.



. Д. Соломенцев.



ФУ РЬЕ - БЕССЕЛЯ РЯД - разложение функции $f(x)$ в ряд



\[

f(x)=\sum_{m=1}^{\infty} c_{m} J_{v}\left(x_{m}^{(v)} \cdot \frac{x}{a}\right), 0<x<a,

\]



где $f(x)$-заданвая в интервале $(0, a)$ функция, $J_{v}(x)-$ Бесселя функция порядка $v>-\frac{1}{2}, x_{m}^{(v)}$-положительные нули функции $J_{v}(x)$, расположенные в порядке возрастания; коэффициенты ряда $c_{m}$ имеют следующие значения



\[

c_{m}=\frac{2}{a^{2} J_{v+1}^{2}\left(x_{m}^{(v)}\right)} \int_{0}^{a} r f(r) J_{v}\left(x_{m}^{(v)} \cdot \frac{r}{a}\right) d r .

\]



Если $f(x)-$ кусочно вепрерывная функция, заданная на интервале $(o, a)$ и интеграл



\[

\int_{0}^{a} \sqrt{r}|f(r)| d r<\infty,

\]



то Ф.- Б. р. сходится и сумма его равна $\frac{1}{2}(f(x+0)+$ $+f(x-0))$ в каждой внутренней точке $x$ интервала $(0, a)$, в окрестности к-рой $f(x)$ имеет ограниченную вариацию. Л. Н. кармазина.



ФУРЬЕ - СТИЛТЬЕСА ПРЕОБРАЗОВАНИЕ - ОДно из интегральных преобразований, родственное Фурье преобразованию. Пусть функция $F(x)$ имеет ограниченное изменение на $(-\infty,+\infty)$. Функция



\[

\varphi(x)=\frac{1}{\sqrt{2 \pi}} \int_{-\infty}^{+\infty} e^{-i x y} d F(y)

\]



$\triangle 24$ Математическая әнц., т. 5





\end{document}